\documentclass[10pt,a4paper]{amsart}
\usepackage[margin=19mm]{geometry}
\usepackage{amsmath,amssymb,mathtools}
\usepackage[scr=rsfs,cal=euler]{mathalfa}
\usepackage{xfrac}
\usepackage[unicode,hidelinks,bookmarks=false]{hyperref}
\newcommand{\ie}{i.e.\ }
\newcommand{\cf}{cf.\ }
\newcommand{\resp}{respectively\ }
\newcommand{\Subsection}{\subsection}
\providecommand{\nobreakdash}{\nobreak\hbox{-}}
\setlength{\emergencystretch}{2em}
\multlinegap0pt
\usepackage{amssymb}
\newtheorem{prop}{Proposition}
\newcommand{\cO}{\mathcal O}
\title{Nearly Gorenstein normal graded rings}
\author{Tomohiro Okuma, Kei-ichi Watanabe, Ken-ichi Yoshida}
\date{Source: arXiv:2602.04222v2 (CC BY 4.0)}
\begin{document}
\maketitle

\noindent\emph{Source and licence.} Nearly Gorenstein normal graded rings. Version arXiv:2602.04222v2 (CC BY 4.0), distributed by arXiv under the Creative Commons Attribution 4.0 International licence, http://creativecommons.org/licenses/by/4.0/, as recorded in the arXiv licence metadata for this exact version and shown on the abstract page https://arxiv.org/abs/2602.04222v2. Full licence notice: \texttt{licenses/arxiv-2602.04222-CC-BY-4.0.txt}.\par\smallskip

\noindent\textit{Editorial scope.} Counted unit: the article's prop hypell (source0.tex lines 1176--1205). The article's complete original proof of it is delivered with it (source0.tex lines 1207--1282, one \texttt{proof} environment of 76 lines). No mathematics is written, repaired or re-derived: the statement and the proof are the article's own lines, in the article's own notation, and no retained line is substituted or re-flowed.  No further source-local context is needed: every cross-reference of every retained line resolves inside the two delivered spans.  Nothing was omitted from any retained span, so the package carries no omission record.  No retained line contains a citation, so the wrapper carries no bibliography and no bibliography entry.  No retained line contains a figure environment, an image-inclusion command or drawing code, so no image is delivered and the package declares no asset dependency.  Editorial formatting only, in addition: an AMS \texttt{amsart} preamble that restates the article's own declarations and macros used by the retained lines, namely the environments \texttt{prop} on separate counters, together with the standard packages those lines need.  The retained environments print in this document as Proposition 1 in order of appearance, whereas the article prints the same items with the numbering of its own full document.  The complete native source of the article as distributed in the arXiv e-print for this exact version is bundled unchanged as \texttt{sources/193961/source0.tex}.
\begin{prop}\label{hypell} 
Let $g\ge 2$ and $R= k[X,Y,Z]/(Z^2 + f_{2g+ 1}(X^2,Y))$ with degree 
$\deg (X,Y,Z) = (1,2, 2g+1)$, where $f_{2g+1 }(X^2,Y)$ is a 
homogeneous polynomial of 
degree $2g+1$ of $(X^2, Y)$ with no multiple roots.
Note that $a(R)= 2g-2$ in this case. 
This is a cone singularity $R(C,P)$, where 
$C$ is a hyperelliptic curve of genus $g$ and $P\in C$ is a hyperelliptic point 
$($$\dim_k( H^0(C, \cO_C(2P))=2)$. 
Note that $R^{(d)}$ is Gorenstein if and only if $d \,|\, (2g-2)$.   
\par
In the following, we always assume that $d$ does not divide $2g-2$. 
\begin{enumerate}
\item[(1)] If $d=3$ or $4$, then $R^{(3)}, R^{(4)}$ is nearly Gorenstein for every $g$.
\item[(2)] If $d=6$, then $R^{(6)}$ is nearly Gorenstein 
if and only if $g \not\equiv 2 \pmod{3}$. 
\item[(3)] If $d\ge 4g-1 = \deg K_C + \deg Z$, then $R^{(d)}$ is nearly Gorenstein. 
\item[(4)] In the following, we assume that $d\ge 5$. 
\begin{enumerate}
\item[{\rm (a)}] If $g=2$, then $R^{(d)}$ is nearly Gorenstein if and only if $d\ge 7$.
\item[{\rm (b)}]  If $g=3$, then $R^{(d)}$ is nearly Gorenstein if and only if $d=5, 6$ or 
$d\ge 11$.
\item[{\rm (c)}] 
 If $g=4$, then $R^{(d)}$ is nearly Gorenstein if and only if 
$d =5,7, 8$ or $d\ge 15$.
\item[{\rm (d)}]  If $g=5$, then $R^{(d)}$ is nearly Gorenstein if and only if 
$d \le 10, d\ne 6$ or $d\ge 19$.
\end{enumerate}
\end{enumerate}
\end{prop}

\begin{proof} 
Fix $d$ and put $L=K_R^{-1}$. 
\par
Let $K^{(d)}_n$ be the degree $n$ part of 
$K_{R^{(d)}}= [K_R]^{(d)}$, that is, 
\[
K^{(d)}_n=H^0(C, \cO_C(K_C + (nd)P)) = 
H^0(C, \cO_C((2g-2+ nd)P)). 
\]
Let  $L^{(d)}_n$  be the degree $n$ part of 
$K^{-1}_{R^{(d)}}$, that is, 
\[
L^{(d)}_n=H^0(C, \cO_C(-K_C + (nd)P)) = 
H^0(C, \cO_C((-2g+2+ nd)P)). 
\]

\par
(1) First suppose $d=3$. 
If $3 \not | \; 2g-2$, we can find $n$ such that $(2g-2) - 3n = 1$ or $2$, so that
$ K^{(3)}_{-n} \cdot L^{(3)}_{n+1} = R_3$ and similarly for $R_6$. 
\par
Note that if  $2g+1 = \deg Z$ is divisible by $3$, then so is $2g-2$ and hence 
$R^{(3)}$ is Gorenstein. Then it is easy to see that for every $m$, 
we can find $n$ so that $R_{3m} = K^{(3)}_n\cdot  L^{(3)}_{m-n}$. 
\par \vspace{1mm}
Next suppose $d=4$. If $g$ is odd, since $2g-2$ can be 
divided by $4$, then $R^{(d)}$ is Gorenstein. 
If $g=2m$ is even, then we can see $K^{(4)}_{-m+1}\cdot L^{(4)}_{m} = (R_2)^2 =R_4$.
Next, $m+1$ is the smallest degree $n$ such that $Z$ appears in 
$ R^{(4)}_{n} = R_{4n}$ and 
we can see that $K^{(4)}_{1}\cdot L^{(4)}_{m}= R^{(4)}_{m+1}$.  

\par \vspace{1mm}
(2) If $g \equiv 2 \pmod{3}$, then $2g + 1 = 6m -1$ for some positive integer $m$ 
and $XZ \in R^{(6)}_m = R_{6m}$.  But if $s= \pm (2g-2) + 6n < 6m$, then $R_s$ does not 
contain $Z$ and hence $XZ$ does not appear in $K^{(6)} \cdot  L^{(6)}$ and so $R^{(6)}$ is not nearly Gorenstein.  

\par \vspace{1mm}
(3) If $d \ge 4g-1 = \deg Z + \deg K_C$, then it is easy to see that 
$R_d = K^{(d)}_{0}\cdot L^{(d)}_{1}$. Since   $R^{(d)}$ is generated by $R_d$, 
$R^{(d)}$ is nearly Gorenstein.  
\par 
Note that if $g \equiv 1 \pmod{3}$ then $6 \,|\, 2g-2$.  
\par \vspace{1mm}
(4) (a) If $g =2$ and $d=5$, then $Z \in R_5$ is not contained in 
$K^{(5)} \cdot  L^{(5)}$
and $R^{(5)}$ is not nearly Gorenstein. 
Similarly, $R^{(6)}$ is not nearly Gorenstein because 
$ZX \in R_6 \setminus  K^{(6)} \cdot  L^{(6)}$. 
If $d\ge 7$,  then the assertion follows from (3). 

\par 
(b) If $g=3$, then $\deg K_C =4$ and $\deg Z =7$. If $d=5$, then 
$R_5 = K^{(5)}_0 \cdot  L^{(5)}_1$ and 
$R_{10} = K^{(5)}_1 \cdot  L^{(5)}_1 + 
K^{(5)}_0 \cdot  L^{(5)}_2$ imply that $R^{(5)}$ is nearly Gorenstein. 
We can show that $R^{(7)}$ (resp. $R^{(8)}$) is not nearly Gorenstein, 
since $Z \in R_7$ (resp. $XZ \in R_8$) is not contained in 
$K^{(8)} \cdot  L^{(8)}$. 
(resp. $K^{(8)} \cdot  L^{(8)}$). $R^{(9)}, R^{(10)}$ is not nearly Gorenstein in the same manner. 
\par  
(c) If $g=4$, then  $R$ is generated by $X,Y,Z$ with degrees $1,2, 9$ and $\deg (K_C) = 6$.  
We consider the cases $d = 5,7,8, \ldots$.
\par
If $d = 4,5,7,8$, we can show that $R^{(d)}$ is nearly Gorenstein. 
If $9\le d \le 14$, we put  
$d = 9+s$. 
Then we can easily see that $X^s Z\in R_d$ is not in  $K^{(d)} \cdot  L^{(d)}$.
Hence  $R^{(d)}$ is not nearly Gorenstein. If $d \ge 15$, then $R^{(d)}$ is nearly Gorenstein
by (3). 
\par 
% KY0915
%(d) If $g=5$, then then  $R$ is generated by $X,Y,Z$ with degrees $1,2, 11$ and $\deg (K_C) = 8$.
(d) If $g=5$, then  $R$ is generated by $X,Y,Z$ with degrees $1,2, 11$ and $\deg (K_C) = 8$.
We can have the result by the same argument as in (a), (b), (c). 
\end{proof} 

\end{document}
